\section{Entry into the bounded-quotient regime}

\begin{lemma}[Entry]
For every $m\ge1$, there exists
\[
 n_0\le \lceil\sqrt{2m}\,\rceil+2
\]
such that $b_{n_0}<n_0^2$. Once $b_n<n^2$ holds, it continues to hold at every later index.
\end{lemma}

\begin{proof}
Let $F(n)=b_n-n^2$. Since $r_n\le n-1$,
\[
 F(n+1)-F(n)=r_n-(2n+1)\le-(n+2)<0.
\]
Thus $F$ decreases strictly. Starting from $F(2)=m-4$, summing the displayed decrease gives a negative value within at most $\lceil\sqrt{2m}\rceil$ further steps. Strict decrease gives forward invariance.
\end{proof}

\begin{lemma}[Bounded quotient and three possible changes]
If $q_n\le n$, then
\[
 \Delta q_n:=q_{n+1}-q_n\in\{-1,0,1\},
 \qquad q_{n+1}\le n+1.
\]
Consequently, the region $q_n\le n$ is forward invariant, and every orbit eventually enters it.
\end{lemma}

\begin{proof}
Set $d=2r_n-q_n$. If $q_n\le n$, then
\[
 -(n+1)<d<2(n+1).
\]
The quotient correction in \eqref{eq:transition} is therefore one of $-1,0,1$, and $q_{n+1}\le q_n+1\le n+1$.

For the last sentence of the statement, the entry lemma supplies an index with $b_n<n^2$, and then $q_n=\lfloor b_n/n\rfloor<n$; forward invariance of $b_n<n^2$ therefore places every orbit permanently in the region $q_n\le n$.
\end{proof}

The coordinate \eqref{eq:e} satisfies an exact doubling law.

\begin{proposition}[Exact doubling coordinate]
After entry into $q_n\le n$,
\begin{equation}
 e_{n+1}=2e_n-\Delta q_n(n+2),
 \qquad
 e_{n+1}\equiv2e_n\pmod{n+2}.
 \label{eq:doubling}
\end{equation}
\end{proposition}

\begin{proof}
Since $2r_n-q_n=q_n+2e_n$, equation \eqref{eq:transition} gives
\[
 r_{n+1}=q_n+2e_n-\Delta q_n(n+1),
 \qquad q_{n+1}=q_n+\Delta q_n.
\]
Subtracting proves \eqref{eq:doubling}.
\end{proof}
